%---------------------------------------------------------------------------%
%->> Backmatter
%---------------------------------------------------------------------------%
\chapter[致谢]{致\quad 谢}\chaptermark{致\quad 谢}% syntax: \chapter[目录]{标题}\chaptermark{页眉}
%\thispagestyle{noheaderstyle}% 如果需要移除当前页的页眉
%\pagestyle{noheaderstyle}% 如果需要移除整章的页眉

首先，致以最诚挚的谢意给我的导师——杨旭三特聘研究员。是杨老师领我走进了光学显微成像的大门，三年来，他严谨求真的治学态度与敏锐创新的科研思维，始终是我前行路上的榜样。杨老师具备极高的学术洞察力，无论是光学、电学还是生物实验，总能一针见血地指出症结所在；他对前沿技术的开放态度以及对项目可行性的精准把控，让我在此期间的科学训练与逻辑培养中获益良多。不仅在学术上倾囊相授，杨老师在生活中的悉心关怀与为人处世上的言传身教，更让我在个人品格上不断成长。感谢杨老师一直以来的支持、鼓励与殷切期盼，这是我走到今天的坚实后盾。

感谢北京师范大学的李明老师，在我研三这一年人生的重要时期，不仅为我提供了优良的实验环境、安静而独立空间，更给予了我极大的支持与开导。与李老师的每一次探讨，都加深了我对神经科学的理解并激发了浓厚的兴趣。您对科研的纯粹热爱与对未知事物的不懈探索深深打动了我，同时您的处世准则和对学生的指导完美诠释了``学为人师，行为世范''的北师大校训精神。同时，衷心感谢中国科学院物理研究所的谢浩老师、南京航空航天大学的王德鹏老师，您们上课时候的谆谆教诲与耐心指导让我受益颇多。以及，诚挚感谢哥廷根大学(Georg-August-Universität Göttingen)的 Prof. Dr. Jörg Enderlein。在我最初探索荧光寿命成像技术而感到迷茫时，是您分享的第一份示例代码为我的研究扫清了初期的障碍。您在显微成像领域的深厚造诣与无私分享的学者风范深深启发了我，让我更加坚定地在这一领域继续探索。

感谢长三角物理研究中心的行政助理彭春蕾姐姐，三年来在工作流程和日常生活上给予了我莫大的帮助，在为人处事上教会了我非常多。感谢西湖大学的宋云峰师兄，您时常的关心与对我个人成长的殷切嘱托令我倍感温暖，在此祝愿师兄前程似锦。感谢刘锴老师、秦怀远、蔡江山、柴付美、马龙洋、董钠、管伟楠、李东润、李博等各位老师，感谢你们在专业领域提供的渊博知识、技术支持以及有益讨论。感谢中国科学院物理研究所的贺一诚师兄在超快激光压缩及光路搭建的指导。感谢同门师兄弟王豪、方祎航、李小凡、曾好、王泽宇等的互助与关照，与你们并肩作战的时光不仅开阔了我的视野，更让我的硕士生涯充满了快乐与美好的回忆。

感谢北京师范大学的刘诗语学姐、丁洪岩、赵颖雯、王潇等同学以及各位师弟妹们，在北师大交流的这几个月里多有打扰，感谢你们的热情支持与帮助。感谢井瑞师兄时常的开导与交流，您渊博的学识和宽广的视野为我指点了迷津，提供了巨大的帮助。

饮水思源，感谢中国科学院物理研究所和长三角物理研究中心为我提供的卓越平台与各项支持，让我能够心无旁骛地潜心实验、度过充实快乐的三年时光。特别感谢汪庭语老师、唐庆老师、陈冶蕲老师、秦晓宇老师、高子尚老师、张亚丽老师等各位老师，感谢你们在学业与生活方方面面给予的指导与关怀。

最后，我要将最深的感激献给我的父母。感谢你们多年来的关爱与包容，让我能够无忧无虑地追逐自己的梦想。是你们在我受挫时给予指引，在遇阻时倾力相助，在徘徊迷茫时提供最温暖的避风港。正是你们的无私付出与默默支持，赋予了我在求学之路上披荆斩棘、勇往直前的力量，让我能够自信地迈向下一段科研旅程。


\rightline{2026年3月 于北京师范大学}
\rightline{胡子睿}


\chapter{作者简历及攻读学位期间发表的学术论文与其他相关学术成果}

\section*{作者简历:}
2019年09月——2023年06月，在北京航空航天大学物理学院获得理学学士学位

2023年09月——2026年06月，在中国科学院物理研究所攻读电子信息硕士学位

% \section*{已发表(或正式接受)的学术论文:}

% {
% \setlist[enumerate]{}% restore default behavior
% \begin{enumerate}[nosep]
%     \item 已发表的工作1
%     \item 已发表的工作2
% \end{enumerate}
% }
\section*{申请或已获得的专利:}
\begin{enumerate}[nosep]
    \item 杨旭三, \textbf{胡子睿}。 一种优化多模态成像的方法及多模态成像系统。中国专利, 公开号：CN119048460A
    \item 杨旭三, \textbf{胡子睿}, 等。 基于光声信号辅助自适应光学优化多模态成像的方法。中国专利, 公开号：CN119033299A
    \item 杨旭三, \textbf{胡子睿}, 等。 一种基于啁啾脉冲放大的四模态同步成像系统及方法。中国专利, 公开号：CN120213922A
    \item 杨旭三, \textbf{胡子睿}, 等。 一种提高光学神经网络细胞筛选装置光子利用效率的方法。中国专利, 公开号：CN120064286A
    \item 杨旭三, \textbf{胡子睿}, 等。 一种基于光强选通的腔道内高动态范围内窥成像装置。中国专利, 公开号：CN120314215A
    \item 杨旭三, 王豪, \textbf{胡子睿}。 一种光谱仪专用的暗箱出线孔遮光工装。中国专利, 公开号：CN223692234U
\end{enumerate}

\section*{参加的研究项目及获奖情况:}
\begin{enumerate}[nosep]
    \item 2023-2024 中国科学院大学 优秀学生干部
    \item 2024-2025 中国科学院大学 优秀学生干部
    \item 2025年 中国科学院物理研究所 所长奖学金
\end{enumerate}


\cleardoublepage[plain]% 让文档总是结束于偶数页，可根据需要设定页眉页脚样式，如 [noheaderstyle]
%---------------------------------------------------------------------------%